\section*{الفصل \ref{ch:g7:prop} --- التناسب}

\begin{solution}{exo:g7:prop:1}
الجدول الأول: \omterm{def:g3:division:remainder}{خوارج القسمة} $\frac{7.5}{3} = 2.5$،
و $\frac{17.5}{7} = 2.5$، و $\frac{27.5}{11} = 2.5$: إذن متناسب،
ومعامله $2.5$.

الجدول الثاني: $\frac62 = 3$ و $\frac{15}{5} = 3$، لكن
$\frac{28}{9} \approx 3.1$: إذن غير متناسب.
\end{solution}

\begin{solution}{exo:g7:prop:2}
بقاعدة الضرب المتقاطع: $5 \times p = 6 \times 8$، إذن
$p = \frac{48}{5} = 9.60$ يورو.
\end{solution}

\begin{solution}{exo:g7:prop:3}
الوتيرة: $36 \div 3 = 12$ صفحة في الدقيقة. وفي $5$ دقائق:
$60$ صفحة. ومن أجل $96$ صفحة: $96 \div 12 = 8$ دقائق.
\end{solution}

\begin{solution}{exo:g7:prop:4}
$30\,\%$ من $250$: $0.3 \times 250 = 75$. \quad
و $15\,\%$ من $60$: $0.15 \times 60 = 9$. \quad
و $t\,\%$ من $80$ هي $20$: $\frac{20}{80} = \frac{25}{100}$، إذن
$t = 25$.
\end{solution}

\begin{solution}{exo:g7:prop:5}
$\frac{18}{25} = \frac{18 \times 4}{25 \times 4} = \frac{72}{100} =
72\,\%$.
\end{solution}

\begin{solution}{exo:g7:prop:6}
المسار: $9 \times 25\,000 = 225\,000$ سم $= 2.25$ كم.
والبحيرة: $2$ كم $= 200\,000$ سم على الأرض، إذن
$200\,000 \div 25\,000 = 8$ سم على الخريطة.
\end{solution}

\begin{solution}{exo:g7:prop:7}
$4.3$ م $= 430$ سم، و $430 \div 43 = 10$ سم.
\end{solution}

\begin{solution}{exo:g7:prop:8}
الشوكولاتة لكل حصة: $240 \div 8 = 30$ غ. ومع $300$ غ:
$300 \div 30 = 10$ حصص بالضبط.
\end{solution}

\begin{solution}{exo:g7:prop:9}
\emph{1.} التخفيض: $20$ يورو، أي\ $\frac{20}{80} = 25\,\%$ من الثمن
الأصلي.

\emph{2.} الارتفاع البالغ $20$ يورو يُقارن الآن بالمرجع
\emph{الجديد} $60$: $\frac{20}{60} = \frac13 \approx 33\,\%$. فالنسبة المئوية
تُنسب دائمًا إلى كلّ؛ وقد تغيّر الكلّ، فتغيّرت النسبة
أيضًا.
\end{solution}

\begin{solution}{exo:g7:prop:10}
المعامل من العمود الكامل: $\frac{17.5}{10} = 1.75$. ثم
$y = 4 \times 1.75 = 7$ و $x = \frac{10.5}{1.75} = 6$.
\end{solution}

\begin{solution}{exo:g7:prop:11}
\emph{1.} بعد نسخة واحدة: $10 \times 0.8 = 8$ سم. وبعد نسختين:
$8 \times 0.8 = 6.4$ سم.

\emph{2.} التصغير الثاني ينطبق على النسخة المصغَّرة أصلًا،
لا على الأصل: فالمعاملان يتضاربان،
$0.8 \times 0.8 = 0.64$، إذن نسختان تصغّران إلى $64\,\%$ من
الأصل --- لا $60\,\%$.
\end{solution}

\begin{solution}{pb:g7:prop:1}
\textbf{1.} الارتفاعات والظلال متناسبة، إذن بقاعدة
الضرب المتقاطع (\cref{thm:g7:prop:cross})، مع عمود العصا
$(1,\ 0.4)$ وعمود الشجرة $(h,\ 3.2)$:
$0.4 \times h = 1 \times 3.2$، ومنه
$h = \frac{3.2}{0.4} = 8$ م.

\textbf{2.} البرج: $\frac{26}{0.4} = 65$ م. وظلّ الطفل:
$1.5 \times 0.4 = 0.6$ م (فالظلّ $=$ الارتفاع $\times$ معامل
اللحظة $0.4$).

\textbf{3.} كلما تحرّكت الشمس في السماء تغيّر المعامل
الذي يربط الارتفاعات بالظلال؛ ولا تشترك في المعامل نفسه إلا
القياسات المأخوذة في اللحظة نفسها.

\textbf{4.} في تلك اللحظة يكون المعامل $1$: فكل اِرتفاع
يساوي ظلّه. ومنه يكون اِرتفاع قمّة \omterm{def:g5:solids:def}{الهرم}
$147$ م --- وهو اِرتفاع \omterm{def:g5:solids:def}{الهرم} الأكبر (ويكفي طرف ظلّه
مقيسًا من النقطة الواقعة تحت القمّة).

\textbf{5.} \omterm{def:g3:division:remainder}{خارج قسمة} العمود
$\frac{\text{الظلّ}}{\text{الارتفاع}}$ --- وهو معامل \omterm{def:g7:prop:table}{تناسب}
اللحظة --- مشترك بين كل شيء
عمودي: وذلك بالضبط ما يجعل جدول الارتفاعات
والظلال \omterm{def:g7:prop:table}{جدول تناسب} (\cref{def:g7:prop:table}).

\textbf{6.} تصل أشعة الشمس متوازية. وعلى أرض \emph{مسطّحة}،
كان شعاعان متوازيان يصطدمان بعصوين عموديتين بالزاوية
نفسها: فتلقيان ظلّين \omterm{def:g6:propdata:def}{متناسبين} --- إما لا ظلّ
لكلتيهما، وإما ظلّ لكلتيهما. أما عصا بلا ظلّ (أسوان)
وعصا بظلّ (الإسكندرية)، في اللحظة نفسها، فأمر
مستحيل على أرض مسطّحة: فلا بدّ أن يشير العموديان في
اِتجاهين مختلفين --- أي إن السطح ينحني.

\textbf{7.} على الرسم، يشير عمودي أسوان مباشرةً
إلى الشمس؛ وعمودي الإسكندرية مائل بالزاوية
بين اِتجاهَي المدينتين، مرئيةً من المركز. وبما أن
الأشعة متوازية، يكون ذلك الميل بالضبط $7.2^\circ$ المقيسة
بين الشعاع والعصا في الإسكندرية.

\textbf{8.} $\frac{7.2}{360} = \frac{72}{3600} = \frac{1}{50}$:
أي جزء من خمسين من الدورة الكاملة.

\textbf{9.} أطوال الأقواس متناسبة مع الزوايا: فإذا كانت
$7.2^\circ$ --- وهي جزء من خمسين من الدورة --- توافق
$800$ كم، فالدورة كلها توافق
\[
50 \times 800 = 40\,000 \text{ كم}.
\]

\textbf{10.} \omterm{def:g4:geometry:circle}{القطر} $= \frac{\text{المحيط}}{\pi}
\approx \frac{40\,000}{3.14} \approx 12\,700$ كم --- مقابل
$12\,742$ كم الحديثة: أي صحيح إلى ما دون واحد بالمئة بكثير،
بعصا.

\textbf{11.} $40\,000$ كم $= 4\,000\,000\,000$ سم، و
$4\,000\,000\,000 \div 40\,000\,000 = 100$ سم $= 1$ م من
\omterm{def:g6:measure:perimeter}{المحيط}. \omterm{def:g4:geometry:circle}{والقطر}: $100 \div 3.14 \approx 32$ سم --- أي
كرة مكتب أنيقة.

\textbf{12.} $4.8$ كم $= 480\,000$ سم، و
$480\,000 \div 40\,000\,000 = 0.012$ سم $= 0.12$ مم. فأعلى
جبال الألب عُشر ميليمتر على كرة
محيطها متر: أي ذرّة غبار. وعلى هذا \omterm{def:g7:prop:scale}{السلّم} تكون الأرض
أملس من معظم الكرات المصقولة.

\textbf{13.} $40\,000 \div 40 = 1\,000$ يوم، و
$1\,000 \div 365 \approx 2.7$: أي نحو ثلاث سنوات من المشي.

\textbf{14.} $\frac{10}{6\,370} \approx 0.0016$، أي نحو
$0.2\,\%$ (وبدقّة أكبر $0.16\,\%$): فالجوّ قشرة
رقيقة كالهمس على الكوكب.

\textbf{15.} الخطأ: $42\,000 - 40\,000 = 2\,000$ كم، و
$\frac{2\,000}{40\,000} = \frac{5}{100} = 5\,\%$. وفي جملة واحدة:
بعصا وبئر وطريق مقيسة \omterm{def:g7:prop:table}{وتناسب} واحد،
قاس إراتوستينس كوكبًا بدقّة في حدود بضعة بالمئة ---
فالرياضيات تسافر بعيدًا بقليل جدًّا.
\end{solution}
